\section*{Simultaneous shared-coordinate covariance obstruction}

All logarithms are natural.  Let
\[
 q=\|b\|_2^2,\quad Z=\mathcal E_{\rm raw}(b),\quad
 V=\operatorname{Var}(h),\quad J=\mathbf I(h).
\]
The preceding B12 theorem uses the potential
\[
 P=C_*qJ+VZ,\qquad C_*=4+2\log2+2/e,
\]
when source and child coordinates are disjoint.  This note determines what
happens if one physical coordinate is shared.

\subsection*{Exact simultaneous-update formula}

Restrict one shared coordinate and denote its two children by signs \(+,-\).
Use bars and half-differences:
\[
 \bar q=(q_++q_-)/2,\quad \delta q=(q_+-q_-)/2,
\]
and similarly for \(Z,V,J\).  Parseval and the usual restriction identities
give
\[
 q=\bar q,\qquad J=\operatorname{Inf}_i(h)+\bar J,
 \qquad V=\alpha^2+\bar V,
 \qquad Z=\bar Z+D_Z,
\]
where \(D_Z\ge0\) is the scalar-kernel entropy dissipation of \(b\).  Hence,
for arbitrary constants \(A,\lambda\),
\begin{align}
 AqJ+\lambda VZ-\frac{A(q_+J_++q_-J_-)
 +\lambda(V_+Z_++V_-Z_-)}2
 ={}&A\bigl(q\operatorname{Inf}_i(h)-\delta q\,\delta J\bigr)\\
 &+\lambda\bigl(\alpha^2Z+\bar V D_Z-\delta V\,\delta Z\bigr).
 \tag{1}
\end{align}
The last term on each line is a genuine shared-coordinate covariance.  It is
absent when only one block changes.

\subsection*{A four-bit genuine Boolean falsifier}

In lexicographic truth-table order, take
\[
\begin{split}
 g_+={}&(1,-1,-1,1,-1,-1,-1,1,-1,-1,1,-1,-1,-1,-1,1),\\
 g_-={}&(1,-1,-1,-1,1,-1,-1,1,-1,-1,1,-1,-1,-1,1,-1),\\
 h={}&(1,1,1,1,1,1,1,1,1,1,1,1,1,1,-1,1).
\end{split}
\tag{2}
\]
Query coordinate \(i=2\) if coordinates are numbered from one (bit one in
zero-based indexing).  Put \(a=(g_++g_-)/2\),
\(b=(g_+-g_-)/2\), and \(F=a+hb\).

This is not a free-array relaxation.  All three functions are genuine Boolean
signs.  The disagreement support of the two branches is
\(\{3,4,14,15\}\) in zero-based truth-table indexing.  Off that support
\(h=1\), while on it \(h\) takes both signs.  Thus both branches are actually
selected where they differ; the example is not an inactive-child encoding.

Direct rational Walsh calculation gives
\[
 q=\frac14,\qquad Z=\frac54\log2,qquad
 V=\frac{15}{64},\qquad J=\frac12.
\]
The zero-child data are
\[
 (q_0,Z_0,V_0,J_0)=
 \left(\frac18,\frac34\log2,0,0\right),
\]
and the one-child data are
\[
 (q_1,Z_1,V_1,J_1)=
 \left(\frac38,\frac94\log2-\frac9{16}\log3,
       \frac7{16},\frac34\right).
\]
Consequently the two component drops are
\[
 qJ-\frac{q_0J_0+q_1J_1}{2}=-\frac1{64},                         \tag{3}
\]
\[
 VZ-\frac{V_0Z_0+V_1Z_1}{2}
 =\frac{63\log3-102\log2}{512}<0.                                \tag{4}
\]
The sign in (4) is exact because
\(3^{63}<2^{102}\); for instance \(3^5<2^8\) and
\(3^3<2^6\).

Equations (3)--(4) prove:

\paragraph{Proposition (refutation).}
No nonzero potential of the form
\[
 P_{A,\lambda}=Aq\mathbf I(h)+\lambda\operatorname{Var}(h)
                 \mathcal E_{\rm raw}(b),\qquad A,\lambda\ge0,
\]
is a supermartingale under all genuine Boolean shared-coordinate
restrictions.  This includes every attempt to repair the preceding B12
potential merely by increasing the coefficient on its terminal reserve.

The output function itself is harmless:
\[
 H(F)=3\text{ bits},\qquad \mathbf I(F)=2,qquad H(F)/\mathbf I(F)=3/2.
\]
Thus this refutes the two-term potential extension, not FEI.

\subsection*{Unbounded-dimensional tensor lift}

For each \(k\ge0\), let \(s_k\) be the standard bent Boolean sign on
\(2k\) fresh variables,
\[
 s_k(x,y)=(-1)^{x\cdot y}.
\]
It has
\[
 H(s_k)=2k\text{ bits},\qquad \mathbf I(s_k)=k.
\]
Replace \(g_\pm\) by \(g_\pm s_k\) and leave \(h\) unchanged.  Then
\(b\) becomes \(bs_k\), and exact tensor entropy gives
\[
 \mathcal E_{\rm raw}(bs_k)=\mathcal E_{\rm raw}(b)+qH(s_k)\log2.
\]
The two drops become
\[
 \Delta(qJ)=-\frac1{64},
\]
\[
 \Delta(VZ)=
 \frac{63\log3-(102+24k)\log2}{512}<0.                            \tag{5}
\]
Thus the obstruction persists, and the reserve drift tends to
\(-\infty\), in a genuine unbounded-dimensional near-tensor family.  Yet
\[
 \frac{H(Fs_k)}{\mathbf I(Fs_k)}
 =\frac{2k+3}{k+2}\longrightarrow2.
\]
Again this is not an FEI counterfamily.

\subsection*{Consequence for the live proof route}

The disjoint interleaving theorem remains valid: its proof never encounters
the covariance terms in (1).  What is refuted is a direct overlapping-input
extension using only positive multiples of \(qI(h)\) and
\(\operatorname{Var}(h)\mathcal E_{\rm raw}(b)\).  Any repair must retain a
signed joint correction, quotient the composition by representation/activation
equivalence, or charge an actual output split quantity.  A larger coefficient
on the old terminal reserve cannot work.
